\documentclass[a4paper,11pt]{article}
\usepackage[utf8]{inputenc}
\usepackage[T1]{fontenc}
\usepackage[french]{babel}
\usepackage{lmodern}
\usepackage{amsmath,amssymb,amsthm}
\usepackage{mathrsfs}
\usepackage{enumitem}
\usepackage{multicol}

\setlist[enumerate]{label=\textbf{\textcolor{Couleur8}{\arabic*.}}, left=1em}

% Si vous voulez aussi modifier les listes enumerate imbriquées :
\setlist[enumerate,2]{label=\textcolor{Couleur8}{\alph*)}}
\setlist[enumerate,3]{label=\textcolor{Couleur8}{\roman*)}}
\setlist[enumerate,4]{label=\textcolor{Couleur8}{\Alph*)}}
\usepackage{graphicx}
\usepackage{xcolor}
\usepackage{tcolorbox}
\usepackage{geometry}
\usepackage{fancyhdr}
\usepackage{titlesec}
\usepackage{cancel}
\usepackage{ifthen}
\usepackage{tikz}
\usetikzlibrary{arrows.meta}
\usepackage{tkz-tab}
\usepackage{eurosym}

% OPTION PROF/ÉLÈVE
% Décommentez UNE des deux lignes suivantes :
\newboolean{prof}\setboolean{prof}{true}   % Version professeur (avec corrections des devoirs)
%\newboolean{prof}\setboolean{prof}{false}  % Version élève (sans corrections des devoirs)

\definecolor{Couleur1}{HTML}{4A5D7A} %Th
\definecolor{Couleur2}{HTML}{A5B5D6} %Prop
\definecolor{Couleur3}{HTML}{A8C68A} %Méthode
\definecolor{Couleur6}{HTML}{8A9CC1} %Def
\definecolor{Couleur7}{HTML}{E6B459} %Rq Voc
\definecolor{Couleur8}{HTML}{D36740} %Titres
\definecolor{correction}{HTML}{D36740} % Couleur pour les corrections
\definecolor{coopmathscorr}{HTML}{F15929} % Couleur corrections coopmaths

% Configuration des marges
\geometry{
	top=2cm,
	bottom=2cm,
	inner=1.5cm,
	outer=1.5cm,
}

% Police
\renewcommand{\familydefault}{\sfdefault}

% Compteurs
\newcounter{seancenum}
\newcounter{seancenumD}
\setcounter{seancenum}{1}
\setcounter{seancenumD}{1}

% Configuration des en-têtes et pieds de page
\pagestyle{fancy}
\fancyhf{}
\ifthenelse{\boolean{prof}}{
    \fancyhead[L]{\textcolor{Couleur8}{\textbf{Corrections Automatismes - Classe de seconde - Version Professeur}}}
}{
    \fancyhead[L]{\textcolor{Couleur8}{\textbf{Corrections Automatismes - Séance 26 - Classe de seconde}}}
}
\fancyhead[R]{\textcolor{Couleur8}{\textbf{2025-2026}}}
\fancyfoot[L]{\textcolor{Couleur8}{Mme Bertail et Mme Pinto}}
\fancyfoot[R]{\textcolor{Couleur8}{\thepage}}
\renewcommand{\headrulewidth}{1pt}
\renewcommand{\headrule}{\hbox to\headwidth{\color{Couleur8}\leaders\hrule height \headrulewidth\hfill}}

% Environnements personnalisés
\newtcolorbox{seance}[1]{
    colback=Couleur6!10,
    colframe=Couleur1!65,
    fonttitle=\bfseries\large,
    title=Correction Séance \theseancenum #1,
    left=5mm,
    right=5mm,
    top=3mm,
    bottom=3mm,
    arc=0mm,
    after=\stepcounter{seancenum}
}

\newtcolorbox{seanceD}[1]{
    colback=Couleur6!5,
    colframe=Couleur6!45,
    fonttitle=\bfseries\large,
    title=Correction Devoirs - Séance \theseancenumD #1,
    left=5mm,
    right=5mm,
    top=3mm,
    bottom=3mm,
    arc=0mm,
    after=\stepcounter{seancenumD}
}

\begin{document}

\setcounter{seancenum}{26}
\newpage
\begin{seance}{} %26
            \begin{enumerate}[start=19]
	\item $\dfrac{1}{R}=1-\dfrac{2}{2\,026} = \dfrac{1 \times 2\,026}{2\,026} - \dfrac{2}{2\,026} = \dfrac{2\,026}{2\,026} - \dfrac{2}{2\,026}  =\dfrac{2\,024}{2\,026}$. Ainsi, $R={\color[HTML]{f15929}\boldsymbol{\dfrac{2\,026}{2\,024}}}$.
	\item  $\dfrac{2}{2\,026} +\dfrac{1}{4\,052}=\dfrac{4}{4\,052} +\dfrac{1}{4\,052}={\color[HTML]{f15929}\boldsymbol{\dfrac{5}{4\,052}}}$
	\item  On utilise le théorème de Pythagore dans le triangle $ABC$,  rectangle en $B$.\\
              On obtient :\\
              $\begin{aligned}
                AC^2&=BC^2+AB^2\\
                AB^2&=AC^2-BC^2\\
                AB^2&=(\sqrt{2\,026})^2-6^2\\
                AB^2&=2\,026-36\\
                AB^2&=1\,990\\
                AB&= {\color[HTML]{f15929}\boldsymbol{\sqrt{1\,990}}}\\
                \end{aligned}$ 
	\item Pour obtenir $2\,026$, on a ajouté $1\,700$ à $326$. Ensuite, le nombre qui, multiplié par $100$, donne $1\,700$ est $17$.\\
    Le nombre choisi au départ est donc ${\color[HTML]{f15929}\boldsymbol{17}}$.
	\item Une fonction linéaire est une fonction de la forme $f(x)=ax$.\\
    Comme $f(2\,026)=-4\,052$, on a $-4\,052=a\times 2\,026$, soit $a=-2$.\\
    On obtient donc : $f(x)={\color[HTML]{f15929}\boldsymbol{-2x}}$.
	\item La somme de $a$ et $2\,026$ en fonction de $a$ est donnée par ${\color[HTML]{f15929}\boldsymbol{2\,026+a}}$
	\item On obtient l'antécédent de $2026$ en résolvant l'équation $f(x)=2026$.\\
      $\begin{aligned}
      -3x+2\,023&=2\,026\\
      -3x&=3\\
      x&={\color[HTML]{f15929}\boldsymbol{-1}}
      \end{aligned}$
	\item $f(-5)=(-5)^2+2\,026 =25+2\,026 ={\color[HTML]{f15929}\boldsymbol{2\,051}}$
	\item On procède par étapes successives.\\
        On commence par isoler $-3x$ dans le membre de gauche en ajoutant
        $-1\,999$ dans chacun des membres, puis on divise
        par $-3$ pour obtenir la solution : \\
        $\begin{aligned}
        -3x+1\,999&=2\,026\\
       -3x&=2\,026-1\,999\\
       -3x&=27\\
       x&=\dfrac{27}{-3}\\
       x&=-9
       \end{aligned}$\\
       La solution de l'équation $-3x+1\,999=2\,026$ est ${\color[HTML]{f15929}\boldsymbol{-9}}$.
       
	\item  Si $n$ est pair, $(-1)^n=1$ et si $n$ est impair, $(-1)^n=-1$. \\
        Ainsi, $(-1)^{2\,026}+(-1)^{2\,025}=1+(-1)={\color[HTML]{f15929}\boldsymbol{0}}$.
	\item $f({\color[HTML]{f15929}\boldsymbol{2\,025}}) ={\color[HTML]{f15929}\boldsymbol{2\,026}}$
	\item Pour remplir $20$ bouteilles d'huile d'olive, Stéphane utilise $100$ kg d'olives.\\ Cela signifie que pour remplir $1$ bouteille d'huile, il utilise $5$ kg d'olives car $100 \div  20 = 5$.\\On a :\\ 
      $\begin{aligned}
      2\,026&=2\,000+26\\
      &=400\times 5+5\times 5+1\\
      &=405\times 5+1
      \end{aligned}$\\
      Il peut remplir ${\color[HTML]{f15929}\boldsymbol{405}}$ bouteilles d'huile d'olive.
\end{enumerate}
\end{seance}

\end{document}